\chapter{Fazit}
\label{fazit}

Wie im vorigen Kapitel schon erwähnt wurden die Ziele der einheitlichen Anbindung und Ansteuerung heterogener Systeme, prozesseinhaltenden Produktionen und der transparenten Aufbereitung des Produktionsprozesses für alle Stakeholder (siehe Kapitel \ref{motivation-und-zielsetzung}) zum Großteil erreicht. Es wurde aus einem 
Prozessmodell, welches den Soll-Zustand eines Prozesses aufzeigt ein ausführbarer Ablaufplan in Form eines Ausführungsmodells erzeugt. Dieses kann zur Entwicklung von Firmare und 
der daraufhin folgenden Orchestrierung von IoT-Geräten mit der programmierten Firmware verwendet werden. 

\section{Ausblick}

Neben der Umsetzung der Lösungen der kritischen Punkte, gibt es die folgenden möglichen Aussichten für das Resultat dieser Arbeit:

\subsection{Modularisierung durch Einführung von unterschiedlichen Notationsfrontends, Zielbackends und Eventloggingformaten}

Aktuell unterstützt der Protoyp die EPK-Notation, läuft auf ARIS und generiert Code für den ESP32. Jedoch können auch Frontends für andere Notationen 
wie BPMN oder BPMN-2.0 entwickelt werden und auch kann man für unterschiedliche Prozessverwaltungsplattformen das passende Backend entwickeln.
Neben diesen zwei Aspekten kann auch die Codegenerierung für weitere Mikrocontroller wie Arduino oder STM32 ausgeweitet werden. Auch kann das Logging der 
Ausführung in anderen Formaten umgesetzt werden, sodass weitere Auswertung auf Basis dieser Daten durchgeführt werden können.

\subsection{Open-Sourcing des Ausführungsmodells, Codegeneration, Event-Logging und Live-Monitoring}

Der Eigenanteil dieser Arbeit ist der Compiler, das Ausführungsmodell sowie die Codegenerierung und das Live-Monitoring. Dadurch, dass es sehr viel Potenzial gibt das ganze weiter 
auszubauen ist es sinnvoll, diese Anteile (mit Ausnahme der ARIS-Komponenten) der Öffentlichkeit zugänglich zu machen, um die Automatisierung und Orchestierung von industriellen Geschäftsprozessen 
durch Open-Source Entwicklung voranzutreiben.

\section{Praxistauglichkeit}

Durch die Arbeiten in Abschnitt \ref{related-work} hat man schon gesehen, dass Prozesse automatisiert werden und somit die Lücke zwischen Prozessmodellierung und tatsächlicher Ausführung geschlossen werden kann. Jedoch tut diese Arbeit auch neben dieser Prozessautomatisierung, ein Rahmenwerk festlegen, 
sodass Software gemäß von Geschäftsprozessen entwickelt werden kann. Diese prozessorientierte Softwareentwicklung kann in der Praxis sehr sinnvoll sein und durch Vereinfachung von Entwicklungsprozessen und 
Lösungen wie Codegenerierung ist es umsetzbar. Kritisch kann es jedoch werden bei der Aufrüstung der industriellen Komponenten, da hier Anbindungen gemacht werden müssen zwischen der Anlage und dem Mikrocontroller, 
für eine Verbindung zwischen ARIS und den Anlagen wie in Abschnitt \ref{reflection} angesprochen. Es ist also im Fazit umsetzbar softwareseitig, jedoch kann es auf der Hardwareseite zu Schwierigkeiten kommen. Sei es durch Inkompatibilität 
oder einfach der fehlenden Möglichkeit der Anbindung an die Anlage. Es benötigt Anlagen mit Periphärien oder einer Schnittstelle, um einen Mikrocontroller zu integrieren. Hierbei kann man sich ein Beispiel an 
Periphärien wie CAN aus der Automobilindsutrie nehmen, welche zur Realisierung von OBD verwendet wird. \cite{noauthor_road_2024} \cite{hristu-varsakelis_vehicle_2005}
